\subsection{Factoring a Prime Polynomial?}
Back in \unref{UnPolyBasics}, we had we called a polynomial like this one \emph{prime}:
\[
x^2-4x+1
\]
In order to factor this, the \makeblanksmall method says we need to find two numbers whose product is \makeblanksmall and whose sum is \makeblanksmall, and \makeblank, so we \makeblank.

On the other hand, we can use the \makeblank[1in] or \makeblank[1in] to find two perfectly good solutions to the equation
\[
x^2-4x+1=\makeblanksmall
\]
\begin{lpic}{}{3}
\end{lpic}
So according to the factor theorem, we can factor the polynomial as
\[
x^2-4x+1=\Bigl(\makeblank[1.5in]\Bigr)\Bigl(\makeblank[1.5in]\Bigr)
\]
What gives? We need to make a distinction.
\begin{definition}[Prime Polynomial]
We call a polynomial with rational coefficients \term{prime} if it cannot be factored as a product of polynomials with rational coefficients.
\end{definition}
\begin{definition}[Irreducible Polynomial]
We call a polynomial with real coefficients \term{irreducible} if it cannot be factored as a product of polynomials with real coefficients.
\end{definition}
So $x^2-4x+1$ is \makeblank but not \makeblank.